\chapter{Marco Teórico}
\label{ch:marco_teorico}

En este capítulo se hace una revisión de los conceptos teóricos que dan base a esta tesina de grado. Se abordan temas como el diseño de fármacos asistido por computadora, las formas de representar moléculas, las métricas de evaluación y el acoplamiento molecular. También se explica el funcionamiento de los autoencoders variacionales y los métodos de optimización multiobjetivo para buscar soluciones en espacios latentes. Finalmente, se describe el rol de la proteína MRP1 en el cáncer colorrectal y se resume el estado del arte en relación con el diseño computacional de fármacos.

\section{Diseño y Descubrimiento de Fármacos}
\label{sec:diseno_farmacos}

El descubrimiento de nuevos fármacos es un proceso mediante el cual se buscan, identifican, diseñan y validan nuevas moléculas con el objetivo de tratar enfermedades mediante la inhibición o activación de una proteína diana. Se trata de un proceso largo y costoso, que puede requerir entre 10 y 15 años de investigación y desarrollo, con una inversión que puede superar los 1.000 millones de dólares~\cite{hughes2011principles}.

Frente a estos desafíos, en la actualidad el diseño de fármacos involucra con frecuencia el uso de herramientas computacionales. En términos generales, es posible distinguir dos perspectivas complementarias dentro de este campo: por un lado, el pipeline tradicional, basado en etapas sucesivas de identificación y validación de blancos, ensayos y optimización de candidatos; por otro lado, el diseño de fármacos asistido por computadora (CADD), que incorpora técnicas y métodos computacionales y de tecnologías de la información para acelerar y optimizar el proceso~\cite{zhang2025role}. A continuación, se describe el pipeline tradicional (subsección~\ref{subsec:pipeline_tradicional}) y luego se profundiza en el enfoque computacional (subsección~\ref{subsec:cadd}). 

\subsection{Pipeline Tradicional}
\label{subsec:pipeline_tradicional}
Históricamente, el descubrimiento de fármacos tiene su origen en pruebas de ensayo y error sobre extractos naturales o compuestos aislados, evaluados directamente en células o animales, sin un blanco molecular específico previamente definido~\cite{chawla2024computational}.

La figura~\ref{fig:pipeline_tradicional} ilustra este proceso metódico y secuencial surgido con el desarrollo de la industria farmacéutica, el cual en términos generales se agrupa en tres fases: el descubrimiento del fármaco, el desarrollo preclínico y los ensayos clínicos en humanos, hasta alcanzar la aprobación regulatoria necesaria~\cite{biala2023research}.

% =========================================================================
% Figura: Pipeline tradicional de descubrimiento y desarrollo de fármacos
% Inspirado en Chawla et al. (2024), Figura 1.1 y Hughes et al. (2011)
% Paleta de colores accesible Okabe-Ito definida en Configurations/01-Colors.sty
% =========================================================================

\begin{figure}[htbp]
\centering
\resizebox{\linewidth}{!}{%
\begin{tikzpicture}[
    font=\sffamily,
    >=Stealth,
    x=1cm, y=1cm
]

% -------------------------------------------------------------------------
% Parámetros geométricos
% -------------------------------------------------------------------------
\def\w{3.4}      % Ancho de cada bloque
\def\hArrow{1.15}% Altura de la flecha de compuestos
\def\hBadge{0.72}% Altura del distintivo de la etapa
\def\tip{0.45}   % Saliente de la punta de la flecha
\def\wing{0.16}  % Aleta de la flecha

% Coordenadas horizontales
\def\xZero{0.00}
\def\xOne{3.75}
\def\xTwo{7.50}
\def\xThree{11.25}
\def\totalW{14.65}

% =========================================================================
% FILA 1: FLECHAS CON CANTIDAD DE COMPUESTOS (EMBUDO)
% =========================================================================

% --- Etapa 1: 10.000 compuestos (Azul) ---
\filldraw[draw=figBlue, fill=figBlueTint, line width=1.1pt, rounded corners=1pt]
    (\xZero, 0) -- (\xZero + \w - \tip, 0) -- (\xZero + \w - \tip, \wing) -- (\xZero + \w, -\hArrow/2) -- (\xZero + \w - \tip, -\hArrow - \wing) -- (\xZero + \w - \tip, -\hArrow) -- (\xZero, -\hArrow) -- cycle;
\node[align=center, text=darkText] at (\xZero + \w/2 - 0.15, -\hArrow/2) {
    {\fontsize{12}{14}\selectfont \textbf{10.000}}\\[1pt]
    {\footnotesize compuestos iniciales}
};

% --- Etapa 2: 250 compuestos (Naranja) ---
\filldraw[draw=figOrange, fill=figOrangeTint, line width=1.1pt, rounded corners=1pt]
    (\xOne, 0) -- (\xOne + \w - \tip, 0) -- (\xOne + \w - \tip, \wing) -- (\xOne + \w, -\hArrow/2) -- (\xOne + \w - \tip, -\hArrow - \wing) -- (\xOne + \w - \tip, -\hArrow) -- (\xOne, -\hArrow) -- cycle;
\node[align=center, text=darkText] at (\xOne + \w/2 - 0.15, -\hArrow/2) {
    {\fontsize{12}{14}\selectfont \textbf{250}}\\[1pt]
    {\footnotesize candidatos preclínicos}
};

% --- Etapa 3: 5 compuestos (Bermellón) ---
\filldraw[draw=figVermilion, fill=figVermilionTint, line width=1.1pt, rounded corners=1pt]
    (\xTwo, 0) -- (\xTwo + \w - \tip, 0) -- (\xTwo + \w - \tip, \wing) -- (\xTwo + \w, -\hArrow/2) -- (\xTwo + \w - \tip, -\hArrow - \wing) -- (\xTwo + \w - \tip, -\hArrow) -- (\xTwo, -\hArrow) -- cycle;
\node[align=center, text=darkText] at (\xTwo + \w/2 - 0.15, -\hArrow/2) {
    {\fontsize{12}{14}\selectfont \textbf{5}}\\[1pt]
    {\footnotesize moléculas en clínica}
};

% --- Etapa 4: 1 fármaco (Verde azulado) ---
\filldraw[draw=figGreen, fill=figGreenTint, line width=1.1pt, rounded corners=1pt]
    (\xThree, 0) -- (\xThree + \w - \tip, 0) -- (\xThree + \w - \tip, \wing) -- (\xThree + \w, -\hArrow/2) -- (\xThree + \w - \tip, -\hArrow - \wing) -- (\xThree + \w - \tip, -\hArrow) -- (\xThree, -\hArrow) -- cycle;
\node[align=center, text=darkText] at (\xThree + \w/2 - 0.15, -\hArrow/2) {
    {\fontsize{12}{14}\selectfont \textbf{1}}\\[1pt]
    {\footnotesize fármaco aprobado}
};

% =========================================================================
% FILA 2: DISTINTIVOS CON EL NOMBRE DE CADA ETAPA
% =========================================================================

\node[
    draw=figBlue,
    fill=figBlueTint,
    line width=0.9pt,
    rounded corners=3pt,
    minimum width=\w cm,
    minimum height=\hBadge cm,
    text width=\w cm - 0.2cm,
    anchor=north west,
    align=center,
    inner sep=3pt
] at (\xZero, -1.35) {
    \textbf{\small\textcolor{figBlue}{Descubrimiento}}\\[1pt]
    {\footnotesize\textcolor{figBlue}{temprano}}
};

\node[
    draw=figOrange,
    fill=figOrangeTint,
    line width=0.9pt,
    rounded corners=3pt,
    minimum width=\w cm,
    minimum height=\hBadge cm,
    text width=\w cm - 0.2cm,
    anchor=north west,
    align=center,
    inner sep=3pt
] at (\xOne, -1.35) {
    \textbf{\small\textcolor{figOrange!90!black}{Fase preclínica}}\\[1pt]
    {\footnotesize\textcolor{grayText}{(\textit{in vitro} / \textit{in vivo})}}
};

\node[
    draw=figVermilion,
    fill=figVermilionTint,
    line width=0.9pt,
    rounded corners=3pt,
    minimum width=\w cm,
    minimum height=\hBadge cm,
    text width=\w cm - 0.2cm,
    anchor=north west,
    align=center,
    inner sep=3pt
] at (\xTwo, -1.35) {
    \textbf{\small\textcolor{figVermilion}{Fases clínicas}}\\[1pt]
    {\footnotesize\textcolor{figVermilion}{(Fases I--III)}}
};

\node[
    draw=figGreen,
    fill=figGreenTint,
    line width=0.9pt,
    rounded corners=3pt,
    minimum width=\w cm,
    minimum height=\hBadge cm,
    text width=\w cm - 0.2cm,
    anchor=north west,
    align=center,
    inner sep=3pt
] at (\xThree, -1.35) {
    \textbf{\small\textcolor{figGreen!80!black}{Aprobación}}\\[1pt]
    {\footnotesize\textcolor{figGreen!80!black}{regulatoria}}
};

% =========================================================================
% FILA 3: BARRAS GLOBALES DE TIEMPO Y COSTO (MÉTRICAS DEL PROCESO COMPLETO)
% Tonos neutros (slate/gris) para indicar alcance global y evitar confusión
% con los colores de las etapas específicas.
% =========================================================================
\def\hBar{0.75}  % Altura más amplia para dar respiro al texto
\def\tipBar{0.45}% Punta de la flecha
\def\wingBar{0.14}% Aleta de la flecha

% --- Barra 1: Tiempo acumulado (10 a 14 años) ---
\filldraw[draw=grayText, fill=cardBorder!50, line width=1pt, rounded corners=1.5pt]
    (0, -2.65) -- (\totalW - \tipBar, -2.65) -- (\totalW - \tipBar, -2.65 + \wingBar) -- (\totalW, -2.65 - \hBar/2) -- (\totalW - \tipBar, -2.65 - \hBar - \wingBar) -- (\totalW - \tipBar, -2.65 - \hBar) -- (0, -2.65 - \hBar) -- cycle;
\node[text=darkText, align=center] at (\totalW/2 - 0.2, -2.65 - \hBar/2) {
    \textbf{Duración acumulada:} 10 a 14 años de investigación y desarrollo
};

% --- Barra 2: Inversión requerida (> 1.000 millones de USD) ---
\filldraw[draw=grayText, fill=cardBorder!50, line width=1pt, rounded corners=1.5pt]
    (0, -3.75) -- (\totalW - \tipBar, -3.75) -- (\totalW - \tipBar, -3.75 + \wingBar) -- (\totalW, -3.75 - \hBar/2) -- (\totalW - \tipBar, -3.75 - \hBar - \wingBar) -- (\totalW - \tipBar, -3.75 - \hBar) -- (0, -3.75 - \hBar) -- cycle;
\node[text=darkText, align=center] at (\totalW/2 - 0.2, -3.75 - \hBar/2) {
    \textbf{Inversión acumulada:} Más de 1.000 millones de dólares (USD)
};

\end{tikzpicture}%
}
\caption[Pipeline tradicional de diseño y desarrollo de fármacos]{Pipeline tradicional de descubrimiento y desarrollo de fármacos, ilustrando la reducción progresiva de candidatos, tiempos e inversión requerida hasta la aprobación regulatoria. Adaptado de~\cite{chawla2024computational, hughes2011principles}.}
\label{fig:pipeline_tradicional}
\end{figure}


En primera instancia, el descubrimiento temprano de fármacos empieza con la identificación y validación del blanco terapéutico, es decir, la proteína, receptor o vía biológica cuya interacción con el fármaco se espera que tenga un efecto terapéutico. La idea de un blanco puede surgir tanto de la investigación académica y clínica como del sector comercial, y suele requerir años de evidencia de respaldo antes de justificar el inicio de un programa de descubrimiento~\cite{hughes2011principles}.

En esta etapa se diseña un ensayo que mide si un compuesto interactúa con su diana. Este ensayo se aplica sobre grandes bibliotecas de compuestos mediante un cribado de alto rendimiento, un proceso automatizado que permite probar miles de moléculas en poco tiempo. Los compuestos que dan resultado favorable en este cribado son posteriormente evaluados y optimizados, modificando su estructura para mejorar su perfil farmacológico, hasta llegar a una molécula candidata~\cite{hughes2011principles}.

Con la molécula candidata definida, se realiza el desarrollo preclínico, donde se estudia la seguridad y la toxicidad en pruebas \textit{in vitro} y con animales, antes de probarse en seres humanos. Si los resultados son satisfactorios, el compuesto avanza a los ensayos clínicos, evaluando la seguridad y la dosis en un grupo reducido de personas, midiendo su eficacia en pacientes y confirmándola finalmente en poblaciones más grandes~\cite{biala2023research}. Superada esta etapa, el fármaco pasa a la aprobación regulatoria, donde organismos como la ANMAT o la FDA autorizan su comercialización~\cite{biala2023research, duran2021regulatory}.

Este recorrido completo explica por qué el desarrollo de fármacos puede extenderse entre 12 y 15 años y superar los 1.000 millones de dólares de inversión~\cite{hughes2011principles}. Miles de compuestos evaluados en las primeras etapas son rechazados posteriormente por distintos factores: falta de interacción con el objetivo, bajo perfil farmacológico, dificultad de síntesis, entre otros. Solo un número reducido llega a fases preclínicas y de aprobación, lo que evidencia la enorme tasa de descarte de candidatos en cada etapa~\cite{chawla2024computational}.


\subsection{Diseño de Fármacos Asistido por Computadora (CADD)}
\label{subsec:cadd}

El diseño de fármacos asistido por computadora (CADD) se originó en la década de 1960, cuando la capacidad de generar modelos moleculares tridimensionales por medios computacionales permitió superar las limitaciones de las representaciones bidimensionales de los compuestos. Desde entonces, evolucionó hasta dividirse en dos enfoques principales según la información disponible sobre el blanco terapéutico: el diseño basado en la estructura (SBDD), que parte de la estructura tridimensional del blanco, y el diseño basado en el ligando (LBDD), que se basa en un conjunto de compuestos con actividad biológica ya conocida contra ese blanco~\cite{chawla2024computational}.

% =========================================================================
% Figura: Bifurcación de metodologías CADD (SBDD y LBDD) y convergencia en VS
% Inspirado en Chawla et al. (2024) y Macalino et al. (2015)
% Paleta de colores accesible Okabe-Ito definida en Configurations/01-Colors.sty
% =========================================================================

\begin{figure}[htbp]
\centering
\resizebox{0.80\linewidth}{!}{%
\begin{tikzpicture}[
    font=\sffamily,
    >=Stealth,
    x=1cm, y=1cm
]

% -------------------------------------------------------------------------
% 1. NODO RAÍZ SUPERIOR: CADD
% -------------------------------------------------------------------------
\node[
    draw=grayText!80,
    fill=cardBorder!40,
    rounded corners=5pt,
    line width=1.2pt,
    minimum width=4.6cm,
    minimum height=1.0cm,
    align=center,
    text=darkText
] (cadd) at (7.0, 5.95) {
    {\fontsize{12}{14}\selectfont \textbf{CADD}}\\[2pt]
    {\footnotesize\color{grayText} Diseño de Fármacos Asistido por Computadora}
};

% -------------------------------------------------------------------------
% 2. BIFURCACIÓN PRINCIPAL: SBDD Y LBDD
% -------------------------------------------------------------------------
% Rama SBDD (Izquierda - Azul)
\node[
    draw=figBlue,
    fill=figBlueTint,
    rounded corners=4pt,
    line width=1.2pt,
    minimum width=4.8cm,
    minimum height=0.9cm,
    align=center,
    text=darkText
] (sbdd) at (3.6, 4.15) {
    {\fontsize{11}{13}\selectfont \textbf{\color{figBlue} SBDD}}\\[1pt]
    {\footnotesize (\textit{Structure-Based Drug Design})}
};

% Rama LBDD (Derecha - Naranja)
\node[
    draw=figOrange,
    fill=figOrangeTint,
    rounded corners=4pt,
    line width=1.2pt,
    minimum width=4.8cm,
    minimum height=0.9cm,
    align=center,
    text=darkText
] (lbdd) at (10.4, 4.15) {
    {\fontsize{11}{13}\selectfont \textbf{\color{figOrange!90!black} LBDD}}\\[1pt]
    {\footnotesize (\textit{Ligand-Based Drug Design})}
};

% Conexión bifurcada de CADD a SBDD y LBDD con bifurcación equidistante
\path (cadd.south) -- (sbdd.north) coordinate[midway] (cadd_mid);
\coordinate (cadd_split) at (cadd.south |- cadd_mid);
\draw[line width=1.1pt, draw=grayText]
    (cadd.south) -- (cadd_split);
\draw[line width=1.1pt, draw=grayText, -{Stealth[length=5pt, width=4pt]}]
    (cadd_split) -| (sbdd.north);
\draw[line width=1.1pt, draw=grayText, -{Stealth[length=5pt, width=4pt]}]
    (cadd_split) -| (lbdd.north);

% -------------------------------------------------------------------------
% 3. CAJAS DE METODOLOGÍAS / HERRAMIENTAS
% -------------------------------------------------------------------------

% --- Sub-cajas SBDD (2 cajas: Docking y Diseño de novo) ---
\node[
    draw=figBlue!70,
    fill=white,
    rounded corners=3pt,
    line width=0.9pt,
    minimum width=2.1cm,
    minimum height=1.05cm,
    align=center,
    text=darkText,
    font=\footnotesize
] (box_docking) at (2.3, 2.20) {
    \textbf{Docking}\\[2pt]\textbf{molecular}
};

\node[
    draw=figBlue!70,
    fill=white,
    rounded corners=3pt,
    line width=0.9pt,
    minimum width=2.1cm,
    minimum height=1.05cm,
    align=center,
    text=darkText,
    font=\footnotesize
] (box_denovo) at (4.9, 2.20) {
    \textbf{Diseño}\\[2pt]\textbf{\textit{de novo}}
};

% Conexión de SBDD a sus 2 cajas
\path (sbdd.south) -- (box_docking.north) coordinate[pos=0.45] (sbdd_mid);
\coordinate (sbdd_split) at (sbdd.south |- sbdd_mid);
\draw[line width=0.9pt, draw=figBlue!80]
    (sbdd.south) -- (sbdd_split);
\draw[line width=0.9pt, draw=figBlue!80, -{Stealth[length=4.5pt, width=3.5pt]}]
    (sbdd_split) -| (box_docking.north);
\draw[line width=0.9pt, draw=figBlue!80, -{Stealth[length=4.5pt, width=3.5pt]}]
    (sbdd_split) -| (box_denovo.north);

% --- Sub-caja LBDD (1 caja: QSAR) ---
\node[
    draw=figOrange!70,
    fill=white,
    rounded corners=3pt,
    line width=0.9pt,
    minimum width=2.6cm,
    minimum height=1.05cm,
    align=center,
    text=darkText,
    font=\footnotesize
] (box_qsar) at (10.4, 2.20) {
    \textbf{QSAR}
};

% Conexión de LBDD a su caja
\draw[line width=0.9pt, draw=figOrange!80, -{Stealth[length=4.5pt, width=3.5pt]}]
    (lbdd.south) -- (box_qsar.north);

% -------------------------------------------------------------------------
% 4. NODO COMÚN DE CONVERGENCIA: CRIBADO VIRTUAL (VS)
% -------------------------------------------------------------------------
\node[
    draw=figGreen,
    fill=figGreenTint,
    rounded corners=5pt,
    line width=1.2pt,
    minimum width=5.6cm,
    minimum height=0.95cm,
    align=center,
    text=darkText
] (vs) at (7.0, 0.35) {
    {\fontsize{11}{13}\selectfont \textbf{\color{figGreen!80!black} Cribado Virtual (VS)}}
};

% Colectores hacia el Cribado Virtual con altura suficiente para las flechas
\coordinate (vs_in_sbdd) at ($(vs.north)+(-1.6,0)$);
\coordinate (vs_in_lbdd) at ($(vs.north)+(1.6,0)$);
\coordinate (vs_turn_level) at ($(vs.north)+(0, 0.40)$);

% Colector SBDD
\draw[line width=0.9pt, draw=figBlue!70] (box_docking.south) -- ++(0,-0.25) coordinate (b_dock);
\draw[line width=0.9pt, draw=figBlue!70] (box_denovo.south) -- ++(0,-0.25) coordinate (b_deno);
\draw[line width=0.9pt, draw=figBlue!70] (b_dock) -- (b_deno);
\coordinate (b_mid) at ($(b_dock)!0.5!(b_deno)$);
\draw[line width=1.1pt, draw=figBlue, -{Stealth[length=5pt, width=4pt]}]
    (b_mid) -- (b_mid |- vs_turn_level) -- (vs_in_sbdd |- vs_turn_level) -- (vs_in_sbdd);

% Colector LBDD
\draw[line width=1.1pt, draw=figOrange, -{Stealth[length=5pt, width=4pt]}]
    (box_qsar.south) -- (box_qsar.south |- vs_turn_level) -- (vs_in_lbdd |- vs_turn_level) -- (vs_in_lbdd);

\end{tikzpicture}%
}
\caption{Bifurcación de las estrategias de diseño de fármacos asistido por computadora (CADD) en métodos basados en estructura (SBDD) y basados en ligando (LBDD), convergiendo en el cribado virtual. Adaptado de~\cite{chawla2024computational, macalino2015role}.}
\label{fig:bifurcacion_cadd}
\end{figure}


En el diseño basado en la estructura, el blanco se obtiene mediante cristalografía de rayos X, resonancia magnética nuclear (RMN) o, en ausencia de datos experimentales, mediante modelado por homología a partir de una proteína homóloga. Este diseño presenta dos técnicas principales, el diseño de novo, que construye ligandos nuevos dentro de las restricciones del sitio de unión, y el acoplamiento molecular, que evalúa cómo encajan distintos compuestos en ese sitio y califica su afinidad mediante funciones de puntuación~\cite{macalino2015role}. En el diseño basado en el ligando, se asume que las moléculas con estructuras semejantes tienden a exhibir actividades biológicas similares, su metodología más extendida es el modelado QSAR, el cual correlaciona variaciones estructurales con la actividad biológica observada~\cite{chawla2024computational}.

Como se observa en la figura~\ref{fig:bifurcacion_cadd}, el cribado virtual converge con ambos enfoques, ya que puede orientarse a partir de la estructura del blanco o de los ligandos de referencia. Asimismo, este proceso no se restringe a predecir afinidades de unión, sino que cumple un papel esencial al filtrar candidatos según su perfil ADMET (absorción, distribución, metabolismo, excreción y toxicidad), descartando compuestos con baja viabilidad~\cite{macalino2015role}. En años recientes, estas herramientas han incorporado distintas metodologías, tales como redes neuronales profundas para mejorar la capacidad predictiva, así como estrategias híbridas que integran SBDD y LBDD dentro de un mismo flujo de trabajo~\cite{chawla2024computational}.

\section{Representación de Moléculas}
\label{sec:rep_moleculas}

El diseño de fármacos asistido por computadora requiere describir eficientemente la estructura de las moléculas para que los algoritmos puedan procesarlas y analizarlas. En este contexto, existen diversos formatos, que van desde conformaciones tridimensionales y grafos moleculares hasta secuencias lineales. Estas últimas describen la estructura del compuesto mediante una cadena de texto, lo cual facilita su almacenamiento, indexación y procesamiento. Entre las notaciones lineales más extendidas se destacan SMILES y SELFIES, ambas abordadas en este trabajo. 

\subsection{SMILES (Simplified Molecular Input Line Entry System)}
\label{subsec:smiles}

La notación SMILES es una representación lineal ampliamente usada para describir la estructura bidimensional de las moléculas mediante cadenas de texto~\cite{weininger1988smiles}. Cada átomo se representa por su símbolo químico y los enlaces por caracteres especiales. Su sintaxis captura la conectividad molecular, qué átomos están enlazados entre sí, utilizando paréntesis y dígitos para ciclos. Esta notación es concisa, legible y compatible con la mayoría de las bases de datos químicas (PubChem, ChEMBL, ZINC), lo que facilita su integración en flujos de trabajo computacionales.

% =========================================================================
% Figura: Representación de moléculas en SMILES (Aspirina 2D vs. SMILES)
% Paleta de colores accesible Okabe-Ito definida en Configurations/01-Colors.sty
% =========================================================================

\begin{figure}[htbp]
\centering
\resizebox{0.95\linewidth}{!}{%
\begin{tikzpicture}[
    font=\sffamily,
    >=Stealth,
    x=1cm, y=1cm
]

% =========================================================================
% CONTENEDOR PRINCIPAL
% =========================================================================
\draw[draw=cardBorder, fill=white, rounded corners=10pt, line width=0.9pt] 
    (-0.5, -0.2) rectangle (15.8, 6.9);

% =========================================================================
% PANEL IZQUIERDO: ESTRUCTURA 2D (GRAFO MOLECULAR)
% =========================================================================
\node[anchor=north west, font=\bfseries\small, text=darkText] at (-0.1, 6.55) {Grafo 2D};
\node[anchor=north west, font=\scriptsize, text=grayText] at (-0.1, 6.15) {Molécula de aspirina};

% Coordenadas del anillo bencénico (Centro en (2.2, 2.9), R = 1.18 cm)
\coordinate (C1) at (3.22, 3.49); % C1 (ramificación 1)
\coordinate (C6) at (2.20, 4.08);
\coordinate (C5) at (1.18, 3.49);
\coordinate (C4) at (1.18, 2.31);
\coordinate (C3) at (2.20, 1.72);
\coordinate (C2) at (3.22, 2.31); % C2 (ramificación 2)

% Coordenadas ramificación lateral 1 (arriba a la derecha)
\coordinate (O_est)  at (3.98, 4.18);
\coordinate (C_act)  at (4.80, 4.85);
\coordinate (O_carb) at (4.80, 5.92);
\coordinate (C_met)  at (5.85, 4.28);

% Coordenadas ramificación lateral 2 (abajo a la derecha)
\coordinate (C_acd)  at (4.15, 1.76);
\coordinate (O_acd)  at (5.08, 2.30);
\coordinate (OH_acd) at (4.15, 0.62);

% --- HALOS COLOREADOS DETRÁS DE CADA FRAGMENTO ---
% 1. Halo Ramificación lateral (Naranja)
\filldraw[fill=figOrangeTint, draw=figOrange, line width=1.2pt, dashed, rounded corners=8pt]
    (3.55, 3.65) -- (3.55, 6.30) -- (6.45, 6.30) -- (6.45, 3.65) -- cycle;

% 2. Halo Ciclo del grafo (Azul)
\filldraw[fill=figBlueTint, draw=figBlue, line width=1.2pt, dashed, rounded corners=8pt]
    (0.70, 1.25) -- (0.70, 4.55) -- (3.48, 4.55) -- (3.48, 1.25) -- cycle;

% 3. Halo Ramificación terminal (Verde)
\filldraw[fill=figGreenTint, draw=figGreen, line width=1.2pt, dashed, rounded corners=8pt]
    (3.65, 0.15) -- (3.65, 2.55) -- (5.60, 2.55) -- (5.60, 0.15) -- cycle;

% --- ENLACES Y ÁTOMOS 2D ---
% Anillo (Kekulé con enlaces dobles alternados)
\draw[line width=1.5pt, draw=darkText] (C1) -- (C6) -- (C5) -- (C4) -- (C3) -- (C2) -- cycle;
\draw[line width=1.1pt, draw=darkText] ($(C2)!0.15!(C1) + (-0.14, 0)$) -- ($(C1)!0.15!(C2) + (-0.14, 0)$);
\draw[line width=1.1pt, draw=darkText] ($(C5)!0.15!(C6) + (0.07, -0.12)$) -- ($(C6)!0.15!(C5) + (0.07, -0.12)$);
\draw[line width=1.1pt, draw=darkText] ($(C3)!0.15!(C4) + (-0.07, 0.12)$) -- ($(C4)!0.15!(C3) + (-0.07, 0.12)$);

% Índices '1' de cierre de ciclo
\node[font=\fontsize{6.5}{7.5}\selectfont\bfseries, text=white, circle, fill=figBlue, inner sep=1pt] at ($(C1)+(-0.35, 0)$) {1};
\node[font=\fontsize{6.5}{7.5}\selectfont\bfseries, text=white, circle, fill=figBlue, inner sep=1pt] at ($(C2)+(-0.35, 0)$) {1};

% Ramificación lateral 1
\draw[line width=1.5pt, draw=darkText] (C1) -- (O_est) -- (C_act);
\draw[line width=1.3pt, draw=darkText] ($(C_act)+(-0.06,0)$) -- ($(O_carb)+(-0.06,0)$);
\draw[line width=1.3pt, draw=darkText] ($(C_act)+(0.06,0)$) -- ($(O_carb)+(0.06,0)$);
\draw[line width=1.5pt, draw=darkText] (C_act) -- (C_met);

\node[font=\footnotesize\bfseries, text=darkText, fill=figOrangeTint, inner sep=1.5pt, rounded corners=1pt] at (O_est) {O};
\node[font=\footnotesize\bfseries, text=darkText, fill=figOrangeTint, inner sep=1.5pt, rounded corners=1pt] at (O_carb) {O};
\node[font=\footnotesize\bfseries, text=darkText, fill=figOrangeTint, inner sep=1.5pt, rounded corners=1pt] at (C_met) {CH$_3$};

% Ramificación terminal 2
\draw[line width=1.5pt, draw=darkText] (C2) -- (C_acd);
\draw[line width=1.3pt, draw=darkText] ($(C_acd)+(0.05, 0.07)$) -- ($(O_acd)+(0.05, 0.07)$);
\draw[line width=1.3pt, draw=darkText] ($(C_acd)+(-0.05, -0.07)$) -- ($(O_acd)+(-0.05, -0.07)$);
\draw[line width=1.5pt, draw=darkText] (C_acd) -- (OH_acd);

\node[font=\footnotesize\bfseries, text=darkText, fill=figGreenTint, inner sep=1.5pt, rounded corners=1pt] at (O_acd) {O};
\node[font=\footnotesize\bfseries, text=darkText, fill=figGreenTint, inner sep=1.5pt, rounded corners=1pt] at (OH_acd) {OH};

% =========================================================================
% PANEL DERECHO: STRING SMILES (SERIALIZACIÓN LINEAL)
% =========================================================================
\node[anchor=north west, font=\bfseries\small, text=darkText] at (8.5, 6.55) {Secuencia SMILES};
\node[anchor=north west, font=\scriptsize, text=grayText] at (8.5, 6.15) {Serialización lineal del grafo};

% Cápsula SMILES completa destacada
\node[
    draw=cardBorder,
    fill=cardBorder!15,
    rounded corners=6pt,
    line width=0.8pt,
    inner sep=6pt,
    anchor=north west
] (smiles_bar) at (8.5, 5.65) {
    {\fontsize{14}{17}\selectfont\ttfamily
    \colorbox{figOrangeTint}{\textbf{\color{figOrange}CC(=O)O}}%
    \colorbox{figBlueTint}{\textbf{\color{figBlue}c1ccccc1}}%
    \colorbox{figGreenTint}{\textbf{\color{figGreen}C(=O)O}}%
    }
};

% 3 Bloques visuales correspondientes a cada fragmento del grafo
% Bloque 1: Ramificación lateral (Naranja)
\node[
    draw=figOrange,
    fill=figOrangeTint,
    rounded corners=5pt,
    line width=1.1pt,
    anchor=north west,
    minimum width=6.5cm,
    inner sep=5pt
] (box_est) at (8.5, 4.30) {
    \makebox[2.4cm][l]{{\fontsize{12}{14}\selectfont\ttfamily\textbf{\color{figOrange}CC(=O)O}}}%
    {\small\textbf{\color{darkText}Ramificación lateral}}
};

% Bloque 2: Estructura cíclica (Azul)
\node[
    draw=figBlue,
    fill=figBlueTint,
    rounded corners=5pt,
    line width=1.1pt,
    anchor=north west,
    minimum width=6.5cm,
    inner sep=5pt
] (box_rng) at (8.5, 2.95) {
    \makebox[2.4cm][l]{{\fontsize{12}{14}\selectfont\ttfamily\textbf{\color{figBlue}c1ccccc1}}}%
    {\small\textbf{\color{darkText}Estructura cíclica} (\texttt{1...1})}
};

% Bloque 3: Ramificación terminal (Verde)
\node[
    draw=figGreen,
    fill=figGreenTint,
    rounded corners=5pt,
    line width=1.1pt,
    anchor=north west,
    minimum width=6.5cm,
    inner sep=5pt
] (box_acd) at (8.5, 1.60) {
    \makebox[2.4cm][l]{{\fontsize{12}{14}\selectfont\ttfamily\textbf{\color{figGreen}C(=O)O}}}%
    {\small\textbf{\color{darkText}Ramificación terminal}}
};

% =========================================================================
% FLECHAS DE CONEXIÓN ENTRE FRAGMENTO 2D Y BLOQUE SMILES
% =========================================================================
\draw[figOrange, line width=1.5pt, -{Stealth[length=6pt, width=4.5pt]}]
    (6.45, 5.0) to[out=5, in=175] (box_est.west);

\draw[figBlue, line width=1.5pt, -{Stealth[length=6pt, width=4.5pt]}]
    (3.48, 3.15) to[out=10, in=175] (box_rng.west);

\draw[figGreen, line width=1.5pt, -{Stealth[length=6pt, width=4.5pt]}]
    (5.60, 1.45) to[out=-5, in=185] (box_acd.west);

\end{tikzpicture}%
}
\caption{Representación de la molécula de aspirina en notación SMILES y correspondencia con su grafo 2D.}
\label{fig:smiles_aspirina}
\end{figure}


Como se observa en la figura~\ref{fig:smiles_aspirina}, la estructura de la aspirina se divide en tres partes dentro de la secuencia SMILES. A la izquierda se presenta el grafo bidimensional y a la derecha su correspondencia en texto, donde las ramificaciones se indican con paréntesis y el anillo aromático se delimita con números que abren y cierran el ciclo.

SMILES se ha consolidado como la representación estándar en modelos generativos de moléculas basados en secuencias, análogos a los modelos de lenguaje. Arquitecturas basadas en redes neuronales pueden aprender las reglas químicas implícitas en las cadenas SMILES y generar nuevas estructuras válidas~\cite{Kong2022SMILES, Segler2018Generating, Olivecrona2017Molecular}. Se ha demostrado además que el uso de SMILES aleatorizados mejora la diversidad y calidad de las moléculas generadas al ampliar el espacio químico de entrenamiento~\cite{ArusPous2019RandomizedSMILES}.

Sin embargo, SMILES presenta limitaciones. Su sintaxis es sensible a pequeños errores que pueden producir moléculas inválidas o distintas. Además, una misma molécula puede tener múltiples representaciones válidas, lo que complica su comparación directa. 

\subsection{SELFIES (Self-Referencing Embedded Strings)}
\label{subsec:selfies}

SELFIES es una representación lineal alternativa a SMILES, diseñada para ser robusta en términos de validez química. Introducida por Krenn et al. entre 2019 y 2020, su objetivo fue eliminar los errores sintácticos comunes de SMILES. En SELFIES, cada símbolo (token) de la cadena corresponde a un elemento químico o una relación estructural, y la gramática está restringida de manera que cualquier secuencia SELFIES decodifica una molécula válida. Es decir, no existen ``SELFIES inválidos'', y hasta una cadena aleatoria genera una estructura respetando las valencias permitidas~\cite{Krenn2019SELFIES}.

\input{Figures/selfies_aspirina.tex}

Como se observa en la figura~\ref{fig:selfies_aspirina}, la misma molécula se representa mediante una secuencia de símbolos entre corchetes. Cada parte de la estructura se define con instrucciones directas, donde las ramas y el cierre del anillo indican de forma exacta los enlaces necesarios para mantener la validez molecular.

La principal ventaja de SELFIES es esta robustez, lo que resulta crucial en modelos generativos. Si se emplea en un VAE, RNN o GAN, prácticamente todas las moléculas generadas serán químicamente válidas, eliminando la necesidad de filtrar salidas inválidas como ocurre con SMILES~\cite{Krenn2019SELFIES}. Esto permite a los algoritmos enfocarse en optimizar propiedades en lugar de aprender reglas gramaticales. SELFIES puede representar las mismas moléculas que SMILES, evitando combinaciones inválidas.

Entre sus limitaciones, las cadenas SELFIES pueden ser menos intuitivas para químicos, ser más largas que su SMILES equivalente y no reflejar de manera simple ciertas transformaciones moleculares. Al igual que SMILES, no incorpora información 3D ni conformacional, y la validez de la secuencia no garantiza que la molécula sea fácil de sintetizar o realista~\cite{Krenn2019SELFIES}.

\section{Métricas de Evaluación y Propiedades Moleculares}
\label{sec:metricas_evaluacion}

En el diseño de fármacos asistido por computadora, la calidad de las moléculas generadas no depende de una única característica; para que un compuesto sea viable como candidato a fármaco, debe cumplir con una serie de propiedades para asegurar un perfil farmacológico adecuado. Estas propiedades abarcan desde características fisicoquímicas hasta su viabilidad de síntesis y afinidad de unión con el objetivo terapéutico. 

\subsection{Reglas de Lipinski y Propiedades Fisicoquímicas}
\label{subsec:lipinski}
% Regla de los 5 de Lipinski (LogP, Peso Molecular, HBD, HBA, TPSA).

Uno de los primeros criterios para estimar si un compuesto químico puede cumplir alguna función farmacológica o actividad biológica es la Regla de los Cinco (\textit{Rule of Five})~\cite{Lipinski1997Solubility}, propuesta en 1997 a partir de un análisis retrospectivo de las propiedades fisicoquímicas de fármacos que habían alcanzado con éxito ensayos clínicos de fase II. La regla establece que una molécula con buena biodisponibilidad oral debe cumplir con las siguientes condiciones:

\begin{itemize}
    \item Peso Molecular $\leq 500$ g/mol
    \item Coeficiente de partición octanol-agua (LogP) $\leq 5$
    \item Número de aceptores de enlace de hidrógeno $\leq 10$
    \item Número de donadores de enlace de hidrógeno $\leq 5$
\end{itemize}

Esta regla se convirtió en un filtro en las primeras etapas del cribado virtual y diseño de fármacos; sin embargo, no es una regla absoluta, ya que existen fármacos que no cumplen con estas condiciones y son exitosos. Por ello, suele complementarse con métricas más detalladas para evaluar el perfil del compuesto. 

\subsection{QED (Quantitative Estimate of Drug-likeness)}
\label{subsec:qed}
% Cuantificación de la similitud a fármaco.

El Quantitative Estimate of Drug-likeness (QED) es una métrica que permite evaluar la probabilidad de que una molécula sea un fármaco viable a partir del análisis de las propiedades fisicoquímicas de una colección de fármacos aprobados~\cite{Bickerton2012QED}. El QED toma un valor entre 0 y 1, donde un valor cercano a 1 indica una mayor probabilidad de que la molécula sea un fármaco viable.

El QED se construye a partir de ocho propiedades moleculares relevantes en el diseño de fármacos: 

\begin{itemize}
    \item Peso molecular (MW)
    \item Coeficiente de partición octanol-agua (ALOGP)
    \item Número de donadores de enlace de hidrógeno (HBD)
    \item Número de aceptores de enlace de hidrógeno (HBA)
    \item Área de superficie polar (PSA)
    \item Número de enlaces rotables (ROTB)
    \item Número de anillos aromáticos (AROM)
    \item Número de alertas estructurales (ALERTS)
\end{itemize}

Cada una de estas propiedades se transforma mediante una función de deseabilidad, derivada de la distribución de estas propiedades en fármacos aprobados; el QED final se calcula como la media geométrica ponderada de estas ocho funciones de deseabilidad.

QED se ha convertido en una de las métricas más utilizadas en los pipelines de diseño de fármacos, debido a que, a diferencia de los filtros discretos como Lipinski, genera una puntuación suave y continua~\cite{Xia2024MolecularOptimization}.

\subsection{Accesibilidad Sintética (SA Score)}
\label{subsec:sa}
% Facilidad y viabilidad de síntesis en laboratorio.

La accesibilidad sintética es un criterio que evalúa la dificultad de síntesis de una molécula. Esta propiedad constituye un filtro común en el diseño de fármacos asistido por computadora, ya que permite que un compuesto prometedor pueda fabricarse de forma viable y eficiente, afectando directamente tanto la etapa de diseño y optimización de moléculas como la viabilidad de su posterior implementación experimental.

Para cuantificar la accesibilidad sintética se suele utilizar el SA Score~\cite{Ertl2009SAScore}; esta métrica asigna a cada molécula un valor entre 1 (fácil de sintetizar) y 10 (difícil de sintetizar). Dicho valor se calcula a partir del número de fragmentos moleculares y la complejidad de las conexiones entre ellos. En la optimización molecular suele emplearse junto con el QED, donde el QED optimiza la deseabilidad farmacológica y el SA Score la viabilidad práctica de obtenerla~\cite{Xia2024MolecularOptimization}.


\subsection{Fracción de Carbonos \texorpdfstring{$sp^3$}{sp3} (Fsp3)}
\label{subsec:fsp3}
% Grado de saturación tridimensional y complejidad molecular.

\subsection{Métodos de Docking Molecular y Afinidad de Unión}
\label{subsec:docking}
% Fundamentos de acoplamiento molecular y cálculo de energía libre de unión con AutoDock Vina.

\section{Modelos Generativos Aplicados al Diseño de Fármacos}
\label{sec:modelos_generativos}

\subsection{Autoencoders Variacionales (VAE)}
\label{subsec:vae}
% Arquitectura encoder-latente-decoder, regularización KL, reconstrucción y muestreo en espacio continuo.

\section{Optimización Multi-Objetivo}
\label{sec:opt_multiobjetivo}

\subsection{Aproximaciones Basadas en la Frontera de Pareto}
\label{subsec:pareto}
% Concepto de dominancia de Pareto, soluciones no dominadas y frentes de Pareto.

\subsection{Indicadores de Rendimiento Multiobjetivo}
\label{subsec:indicadores_multiobjetivo}
% Hipervolumen (HV), Esparcimiento (Spread), Distancia Generacional Invertida (IGD).

\subsection{Metaheurísticas Consideradas en las Experimentaciones}
\label{subsec:metaheuristicas}
% Algoritmos genéticos, PSO, algoritmos evolutivos multiobjetivo aplicados a la búsqueda en espacio latente.

\section{Proteína 1 de Resistencia a Múltiples Fármacos (MRP1)}
\label{sec:mrp1}
% Función biológica de transportadores ABC, mecanismo de expulsión de quimioterápicos en cáncer colorrectal y relevancia de la inhibición de MRP1.

\section{Trabajos Relacionados (Estado del Arte)}
\label{sec:trabajos_relacionados}

\subsection{Diseño de Fármacos Mediante Redes Generativas y VAEs}
\label{subsec:art_vaes}

\subsection{Optimización del Espacio Latente Mediante Metaheurísticas}
\label{subsec:art_metaheuristicas}

\subsection{Diseño de Fármacos Inhibidores de la Proteína MRP1}
\label{subsec:art_mrp1}
